\section{Introduction}
\label{sec:introduction}

Lewy body dementia (LBD) is the clinical umbrella term for two closely related neurodegenerative conditions: dementia with Lewy bodies (DLB) and Parkinson's disease dementia (PDD). Together, they represent the second most prevalent form of neurodegenerative dementia after Alzheimer's disease (AD), accounting for an estimated 15--20\% of all dementia cases at autopsy \cite{mckeith2017dlb}. In the United States alone, LBD affects approximately 1.4 million individuals, while global burden estimates remain substantially higher due to widespread underdiagnosis. The cardinal clinical features of LBD---fluctuating cognition, recurrent visual hallucinations, spontaneous parkinsonism, and REM sleep behaviour disorder (RBD)---reflect the widespread cortical and subcortical distribution of its pathological hallmark: intraneuronal aggregates of the presynaptic chaperone protein alpha-synuclein ($\alpha$-syn), forming structures termed Lewy bodies and Lewy neurites \cite{spillantini1997alpha}.

Despite its prevalence and devastating clinical impact, LBD remains profoundly underdiagnosed and mismanaged. The significant clinical and neuropathological overlap with AD---including the frequent co-occurrence of amyloid plaques and neurofibrillary tau tangles in LBD brains---leads to misdiagnosis rates reportedly between 20 and 80\% in clinical series \cite{mckeith2017dlb}. The consequences of this diagnostic failure are severe: patients with DLB are at risk of life-threatening hypersensitivity reactions to antipsychotic medications commonly prescribed for the behavioural symptoms of AD. Furthermore, no validated blood-based or cerebrospinal fluid (CSF) biomarker currently exists to confirm LBD diagnosis at an early, prodromal stage. The recent addition of $\alpha$-syn seed amplification assays (SAAs) to the biomarker toolkit represents meaningful progress, but SAAs alone do not differentiate LBD subtypes, predict disease trajectory, or identify therapeutic targets.

The molecular aetiology of LBD is multifactorial and complex. At the genetic level, heterozygous loss-of-function mutations in \textit{GBA1} (encoding glucocerebrosidase) are the most common known genetic risk factor, carried by approximately 10\% of patients, while variants in \textit{SNCA}, \textit{APOE} $\varepsilon4$, \textit{TMEM175}, and dozens of additional loci collectively contribute to polygenic susceptibility \cite{chia2021lbd}. A recent large-scale GWAS meta-analysis expanded this risk architecture substantially, identifying 85 LBD risk genes across 51 enriched biological pathways and revealing a novel risk locus, \textit{SYT16}, encoding synaptotagmin-16 \cite{chia2025gwas}. Yet genetic risk accounts for only a fraction of disease variance. Epigenetic modifications---particularly DNA methylation changes at the \textit{SNCA} locus and in neuroinflammatory gene networks---bridge the gap between inherited genetic risk and tissue-level molecular alterations \cite{nabais2024methylation}. Downstream, proteomic and metabolomic profiling in blood and CSF reveals dysregulation in lipid metabolism, lysosomal function, synaptic transmission, and energy metabolism, all of which are pathologically plausible yet poorly understood in the context of LBD progression \cite{frontiers2024metabolomics}.

Multi-omics profiling---the simultaneous collection and integration of data from two or more molecular layers---offers a principled strategy for capturing this multifactorial complexity. Rather than studying each biological stratum in isolation, multi-omics integration identifies convergent molecular signatures that are more robust, more biologically interpretable, and more translationally actionable than single-layer discoveries. Recent computational advances in multi-omics factor analysis (MOFA+) \cite{argelaguet2020mofa}, similarity network fusion (SNF) \cite{wang2014snf}, and deep learning-based embedding methods provide the analytical infrastructure to leverage rich, heterogeneous datasets effectively. Critically, much of the required data already exists in open-access repositories such as the Accelerating Medicines Partnership Parkinson's Disease (AMP-PD) program, the Parkinson's Progression Markers Initiative (PPMI), the UK Biobank, the Gene Expression Omnibus (GEO), and the NIAGADS Genomics Database---substantially lowering the barrier to a purely computational project.

This proposal outlines the scientific rationale, state of the field, proposed methodology, and feasibility assessment for a computational multi-omics project targeting LBD. The remainder of this document is structured as follows: Section~\ref{sec:literature_review} reviews five pivotal recent publications that define the current landscape. Section~\ref{sec:proposed_work} articulates specific research objectives and testable hypotheses. Section~\ref{sec:methodology} describes the computational pipeline in detail. Section~\ref{sec:results} assesses feasibility and projects expected outcomes. Section~\ref{sec:conclusion} provides a synthesis and a recommendation on whether to proceed.
